\subsubsection{Multiscale 3D--1D--0D Coupling}
\label{subsubsec:multiscale-models}

A multiscale model is a coupled forward problem that combines tiers of the
hierarchy. A common arrangement resolves a local 3D stenosis region, connects it
to a surrounding 1D distributed vessel network, and closes terminal beds with
0D Windkessel or related outlet models. The purpose is to keep local resolved
physics where it matters while representing wave propagation and distal
impedance without making the entire circulation a 3D domain
\cite{QuarteroniVeneziani2003GeometricalMultiscale}.

The interface laws are part of the model. A 3D--1D interface may impose
conservation of flow together with pressure, mean-pressure, normal-traction, or
characteristic compatibility. A 1D--0D interface passes terminal flow into an
ODE load and returns a pressure or impedance relation to the 1D boundary.
Resolved-FSI coupling additionally requires compatible wall kinematics,
structural stresses, and interface work. These choices affect stability and
interpretation, and they must be recorded with the numerical coupling convention
\cite{FormaggiaEtAl2007FSICoupling}.

A multiscale model is not a new fidelity tier; it is a coupled forward problem
whose interface laws are part of the model definition. Coupling more components
does not automatically make a comparison more correct. The boundary data,
interface variables, time coupling, output maps, and discrepancy metrics decide
what claim the computation can support. In this report, multiscale coupling
provides comparison and boundary-condition vocabulary; it is not the active
baseline single-vessel solve unless explicitly declared.
